\section{Methodology: Human Evaluation Survey}
\label{ch:methodology}

% ============================================================================
% === WRITING GUIDE: Chapter 4 (~5 pages; prose essentially complete) =======
% ============================================================================
% JOB: argue the measurement gap and present the adapted Direct Assessment
% protocol as the load-bearing instrument. Remaining work is VERIFICATION,
% not writing: (1) the QC-filter subsection now describes the AS-RUN rules
% (completion <10/20, repeat-pair mean |diff| > 25 drop, bad-ref-vs-own-
% real-item-meaning-mean flag); double-check against
% e2r-adaptation/scripts/analyze_survey_responses.py before submission.
% (2) The metaphor-exclusion scoping sentence below restates a supervisor-
% approved decision; the full rationale lives in Sections 2.5 and 7.3.
% THEME HOOK: open: "if no metric can see replacement quality, evaluation
% itself needs explicit structure"; close: "the protocol is deliberately
% conservative; where it limits the claims, the Discussion says so."
% ============================================================================

The survey covers idiom replacement only: the VUAMC provides no gold paraphrases, and the three-part rubric below is well-defined for the figurative-to-literal mapping of idioms but breaks down for conventional metaphors that may require no rewrite at all (Section~\ref{sec:rq3} develops this scoping decision).

This chapter specifies the human-evaluation protocol used to assess figurative-to-literal replacement quality for RQ3 (Section~\ref{sec:rq3}). The replacement task lacks a single ground truth: for an idiom such as \textit{smoking gun}, valid literal paraphrases include \textit{decisive evidence}, \textit{conclusive proof}, or longer rewrites that preserve sentence meaning, none of which need share surface form with each other or with a gold reference. As discussed in Section~\ref{sec:eval-methodology}, automatic metrics in this regime correlate only spuriously with human judgement \cite{Alva-Manchego2021,Scialom2021}, so human evaluation is the primary signal for replacement.

The protocol adopted here is \textit{Direct Assessment} (DA), the standard evaluation method used at the Conference on Machine Translation (WMT), introduced for segment-level machine translation (MT) by Graham et al.\ \cite{Graham2013} and extended to document-level reliability by Graham et al.\ \cite{Graham2017}. The three-part rubric layered on top of DA (grammaticality, meaning preservation, and simplicity) follows the sentence-simplification convention established by Alva-Manchego et al.\ \cite{Alva-Manchego2020}. Section~\ref{sec:da-protocol} describes the protocol; Section~\ref{sec:da-departures} documents three deliberate departures from the original Graham specification, each justified by a difference between the machine-translation setting and the figurative-language replacement setting. The instantiation of the protocol (specific systems compared, source-sentence pool, item counts, and respondent recruitment) is described in the experimental setup of RQ3 (Section~\ref{subsec:rq3-human-eval}).

\subsection{Direct Assessment Protocol}
\label{sec:da-protocol}

\subsubsection{Continuous-scale Rating}

For each evaluated item, the annotator rates the system's literal paraphrase on three independent 0--100 continuous sliders, one per rubric dimension (defined below). The continuous scale is the foundational design choice of the protocol: as Graham et al.\ \cite{Graham2013} demonstrate, replacing a discrete (e.g., 1--5) Likert-style scale with a 0--100 continuous slider unlocks two assessor-intrinsic quality-control techniques that Likert scales cannot support, raises intra-annotator $\kappa$ by up to $+0.101$, and raises inter-annotator $\kappa$ by up to $+0.144$ when combined with per-annotator standardisation. Crowd-sourced annotators on platforms such as Mechanical Turk reach acceptable agreement under this design without expert review.

\subsubsection{The Three Rubric Dimensions}

The three dimensions are rated independently:
\begin{itemize}
    \item \textbf{Grammaticality}: is the output well-formed English?
    \item \textbf{Meaning preservation}: does it convey the meaning of the original sentence?
    \item \textbf{Simplicity}: is the output more accessible than the figurative source?
\end{itemize}
Annotators may optionally provide an alternative paraphrase, retained as qualitative data for the discussion section. Annotators are blind to the system-of-origin; per-condition identifiers live only in analysis-side metadata.

\subsubsection{Standardisation and Aggregation}

Following Graham et al.\ \cite{Graham2013}, raw scores are converted to z-scores per annotator (using each annotator's own mean and standard deviation across all rated items) before aggregation. This corrects for individual scoring tendencies (``score-low'' versus ``score-high'' annotators rating the same outputs) and is the property that makes the continuous scale recover meaningful inter-annotator agreement where Likert-style scales do not.

Aggregation follows the micro-then-macro pattern of \cite{Graham2017}: the mean standardised score is computed per item across its repeat ratings, then per condition across items. Per-item scores are reported as standardised mean $\pm$ 95\% confidence interval, with a composite of the three rubric dimensions as a summary statistic. Confidence intervals use the Student-$t$ distribution over the standardised ratings; given the small per-item repeat counts (Section~\ref{sec:reliability-bound}), they are reported as indicative rather than inferential.

\subsubsection{Quality Control}

Two assessor-intrinsic checks are interleaved into each respondent's set:
\begin{itemize}
    \item \textbf{Bad-reference items}: deliberately degraded versions of system outputs, constructed by lexical substitution or noun-phrase swap that breaks meaning preservation while keeping the sentence superficially fluent. The annotator should rate the bad reference lower than the corresponding genuine system output on the targeted rubric dimension.
    \item \textbf{Repeat-pair items}: duplicate items presented twice in the same respondent's set, separated by other items. The absolute score difference quantifies intra-annotator consistency.
\end{itemize}
Both techniques are introduced in \cite{Graham2013}; the canonical crowd-sourcing implementation (using bad references inside larger HITs, scored by a per-worker difference-of-means significance test, and retaining only workers with $p < 0.05$) is described in \cite{Graham2017}. The thesis uses fixed-threshold filters in place of the per-worker significance test, for reasons documented in Section~\ref{sec:qc-filter-design}.

\subsection{Departures from the Original Protocol}
\label{sec:da-departures}

The DA protocol as described above follows Graham et al.\ \cite{Graham2013,Graham2017} closely. Three design decisions diverge from the original specification, each motivated by a feature of the figurative-language replacement task that does not hold in the machine-translation setting Graham analyses.

\subsubsection{Third Dimension versus Fluency Omission}
\label{sec:rubric-vs-fluency}

Graham et al.\ \cite{Graham2017}, Section 3.3, argue that DA fluency assessment can be safely dropped when resources are constrained: an earlier study \cite{Graham2016a} found no evidence of reference bias in DA adequacy, removing the original methodological motivation for a separate fluency dimension; the freed budget is better spent expanding the test set on the adequacy dimension.

The thesis retains a third dimension (\textit{grammaticality}) in deliberate tension with this argument, on three grounds. First, the task is text simplification, not machine translation: the simplification literature \cite{Alva-Manchego2020} converged independently on a three-axis rubric (grammaticality, meaning preservation, simplicity) because the additional dimension carries information that adequacy alone does not capture. Second, small open-weights LLMs at the scale used here (Llama-3.1-8B and Llama-3.3-70B-Instruct-AWQ; see Section~\ref{subsec:rq3-human-eval}) produce qualitatively different error modes from professional MT systems: agreement and tense errors, dropped articles, and clipped continuations are visible in pilot outputs and would be invisible to a meaning-only rubric. Third, the budget freed by dropping the third axis would translate to perhaps four to five additional items in our 81-item, 22-respondent design, a marginal reliability gain that does not compensate for losing a diagnostic axis.

\subsubsection{Reliability Bound and Exploratory Reporting}
\label{sec:reliability-bound}

Graham et al.\ \cite{Graham2017} establish, by self-replication on Mechanical Turk over the WMT-16 EN--ES document-level QE test set, that DA scores reach $r = 0.901$ between two independent data-collection runs at $\geq 27$ repeat assessments per document (averaging $\approx 107$ segment-level judgements per document). The corresponding segment-level threshold from \cite{Graham2013} is $\geq 15$ repeat assessments. These thresholds define the reliability frontier above which DA can be used as a substitute for an expert gold standard.

The thesis survey ships with a target of $\geq 3$ ratings per item, well below this frontier. This is a deliberate choice driven by the scale of the deliverable: 81 items at 20--25 respondents (16 rated items + 4 quality-control items per respondent) is at the upper end of what a master's-thesis recruitment effort can sustain, and tripling the respondent count to approach Graham's threshold is not feasible. The methodological consequence is that per-condition scores are not used to make reliability-bounded claims about absolute system quality. Instead, the survey is treated as exploratory: results are reported per condition with the standardised means defined above, and pairwise condition comparisons, paired by source sentence, are tested with the paired $t$-test and the Wilcoxon signed-rank test on the per-source mean standardised scores. The survey supports the qualitative claims of RQ2 (decomposition vs.\ monolithic) and RQ1 (within-family scale comparison) at a directional level; it does not establish absolute quality estimates.

\subsubsection{Quality-control Filter Design}
\label{sec:qc-filter-design}

Graham et al.\ \cite{Graham2017} embed bad-reference items inside 100-translation HITs and apply a per-worker difference-of-means significance test, retaining only workers with $p < 0.05$. The test has good power at the original HIT size: with $\sim 100$ items per worker including a balanced bad-reference contingent, a worker who fails to discriminate degraded from genuine outputs is reliably flagged.

The thesis design has 20 items per respondent (16 rated + 4 QC), of which only two are bad references and two are repeat pairs. With $n = 2$ per QC type, the per-respondent significance test has effectively no power: a respondent could rate both bad references higher than their corresponding genuine outputs and still fail to reach $p < 0.05$ under the t-test. Fixed-threshold filters are used instead, applied in three rules. (a)~\textit{Completion}: respondents who fully rate fewer than 10 of their 20 items are excluded. (b)~\textit{Repeat-pair consistency}: respondents whose mean absolute first-versus-second-showing difference across their repeat pairs exceeds 25 points (on the 0--100 scale) are excluded. (c)~\textit{Bad-reference discrimination}: each bad-reference's \textit{meaning preservation} score is compared against the respondent's own mean meaning score across their genuine items; a respondent both of whose bad references score at or above that mean is flagged, and headline results are reported with and without flagged respondents. The bad-reference rule targets the meaning dimension because the degradations are constructed to break meaning preservation specifically while remaining superficially fluent; comparing against the respondent's \textit{own} mean rather than a global threshold preserves the assessor-intrinsic character of the original technique under per-annotator score variation. Threshold values are fixed in advance and reported with the survey; they are not tuned post-hoc against the rated data. The construction recipe for the bad-reference and repeat-pair items is detailed in Appendix~\ref{appendix:qc-construction}.

\newpage
